% !TeX root = ../main.tex
\section{ЛР1 - Перевод параметров из децибел в относительные единицы (разы)}

Для перевода используем формулу $X_{\text{раз}} = 10^{X_{\text{дБ}}/10}$. В расчетах коэффициенты шума ($K_{\text{ш}}$) обозначаются как $F$:
\begin{align*}
\text{Потери фидера: } L_{\Phi} &= 10^{1/10} \approx 1.259 \\
\text{Коэф. усиления преселектора: } K_{p.\text{пр}} &= 10^{2/10} \approx 1.585 \\
\text{Коэф. шума преселектора: } F_{\text{пр}} &= 10^{3/10} \approx 1.995 \\
\text{Коэф. усиления ПЧ: } K_{p.\text{пч}} &= 10^{3/10} \approx 1.995 \\
\text{Коэф. шума ПЧ: } F_{\text{пч}} &= 10^{10/10} = 10 \\
\text{Коэф. шума УПЧ: } F_{\text{упч}} &= 10^{5/10} \approx 3.162 \\
\text{Коэф. усиления УПЧ: } K_{0.\text{упч}} &= 3000 \quad \text{(уже задан в разах)}
\end{align*}

\vspace{1em}
\textbf{2. Расчет коэффициента шума фидера}

Поскольку физическая температура фидера $T_{\Phi} = 293 \text{ К}$ отличается от стандартной шумовой температуры $T_0 = 290 \text{ К}$, коэффициент шума пассивного фидера рассчитывается по точной формуле:
\begin{align*}
F_{\Phi} &= 1 + (L_{\Phi} - 1) \frac{T_{\Phi}}{T_0} \\
F_{\Phi} &= 1 + (1.259 - 1) \frac{293}{290} \approx 1.262
\end{align*}

\vspace{1em}
\textbf{3. Расчет коэффициента шума БВЧ (как системы каскадно-соединенных четырехполюсников)}

Блок высокой частоты (БВЧ) включает в себя фидер, преселектор и преобразователь частоты (ПЧ). Используем формулу Фрииса. Учитывая, что коэффициент передачи фидера $K_{\text{пер.}\Phi} = 1/L_{\Phi}$, формула принимает вид:
\begin{align*}
F_{\text{БВЧ}} &= F_{\Phi} + \frac{F_{\text{пр}} - 1}{K_{\text{пер.}\Phi}} + \frac{F_{\text{пч}} - 1}{K_{\text{пер.}\Phi} \cdot K_{p.\text{пр}}} \\
F_{\text{БВЧ}} &= F_{\Phi} + (F_{\text{пр}} - 1)L_{\Phi} + \frac{(F_{\text{пч}} - 1)L_{\Phi}}{K_{p.\text{пр}}} \\
F_{\text{БВЧ}} &= 1.262 + (1.995 - 1) \cdot 1.259 + \frac{(10 - 1) \cdot 1.259}{1.585} \\
F_{\text{БВЧ}} &= 1.262 + 1.253 + 7.148 \approx 9.663
\end{align*}

\vspace{1em}
\textbf{4. Расчет шумовых температур БВЧ и РПУ}

Эквивалентная шумовая температура БВЧ:
\begin{align*}
T_{\text{БВЧ}} &= T_0 (F_{\text{БВЧ}} - 1) \\
T_{\text{БВЧ}} &= 290 \cdot (9.663 - 1) \approx 2512.3 \text{ К}
\end{align*}

Для расчета суммарной температуры РПУ (которое включает БВЧ и УПЧ), найдем общий коэффициент передачи БВЧ и шумовую температуру УПЧ:
\begin{align*}
K_{\text{БВЧ}} &= \frac{1}{L_{\Phi}} \cdot K_{p.\text{пр}} \cdot K_{p.\text{пч}} = \frac{1}{1.259} \cdot 1.585 \cdot 1.995 \approx 2.512 \\
T_{\text{упч}} &= T_0 (F_{\text{упч}} - 1) = 290 \cdot (3.162 - 1) \approx 627 \text{ К}
\end{align*}

Суммарная шумовая температура приемника (РПУ), приведенная ко входу:
\begin{align*}
T_{\text{РПУ}} &= T_{\text{БВЧ}} + \frac{T_{\text{упч}}}{K_{\text{БВЧ}}} \\
T_{\text{РПУ}} &= 2512.3 + \frac{627}{2.512} = 2512.3 + 249.6 \approx 2761.9 \text{ К}
\end{align*}

\vspace{1em}
\textbf{5. Расчет спектральной плотности мощности приведенного суммарного шума}

Для расчета суммарного шума системы необходимо учесть шумовую температуру антенны $T_A = 300 \text{ К}$. Полная эквивалентная шумовая температура системы:
\begin{align*}
T_{\Sigma} &= T_A + T_{\text{РПУ}} = 300 + 2761.9 = 3061.9 \text{ К}
\end{align*}

Спектральная плотность мощности (СПМ) шума, приведенная ко входу приемника (где $k = 1.38 \cdot 10^{-23} \text{ Дж/К}$ — постоянная Больцмана):
\begin{align*}
N_0 &= k \cdot T_{\Sigma} \\
N_0 &= 1.38 \cdot 10^{-23} \cdot 3061.9 \approx 4.225 \cdot 10^{-20} \text{ Вт/Гц}
\end{align*}

\vspace{1em}
\section{Оценка влияния параметров БВЧ на чувствительность РПУ}

Для повышения чувствительности приемника необходимо уменьшить спектральную плотность мощности (СПМ) приведенного ко входу шума. Исходное значение СПМ: $G_{ш.\Sigma} \approx 4.225 \cdot 10^{-20} \text{ Вт/Гц}$, исходная шумовая температура $T_{\Sigma} = 3061.9 \text{ К}$. 
Целевое значение для пунктов А и Б (уменьшение в 2 раза): 
$G'_{ш.\Sigma} = 2.1125 \cdot 10^{-20} \text{ Вт/Гц}$, что соответствует эквивалентной температуре $T'_{\Sigma} = 1530.95 \text{ К}$. 
Требуемая температура РПУ: $T'_{\text{РПУ}} = 1530.95 - 300 = 1230.95 \text{ К}$.

\textbf{А) Подбор коэффициента передачи преселектора ($K_{p.\text{пр}}$)}

Запишем выражение для $T_{\text{РПУ}}$ в развернутом виде:
\begin{align*}
T_{\text{РПУ}} &= T_0 \left( (F_{\Phi} - 1) + (F_{\text{пр}} - 1)L_{\Phi} + \frac{(F_{\text{пч}} - 1)L_{\Phi}}{K_{p.\text{пр}}} \right) + \frac{T_{\text{упч}} \cdot L_{\Phi}}{K_{p.\text{пр}} \cdot K_{p.\text{пч}}}
\end{align*}
Подставим известные значения в разах ($L_{\Phi}=1.259$, $F_{\text{пр}}=1.995$, $F_{\text{пч}}=10$, $K_{p.\text{пч}}=1.995$, $T_{\text{упч}}=627$):
\begin{align*}
1230.95 &= 290 \left( 0.262 + 0.995 \cdot 1.259 + \frac{9 \cdot 1.259}{K_{p.\text{пр}}} \right) + \frac{627 \cdot 1.259}{K_{p.\text{пр}} \cdot 1.995} \\
1230.95 &= 290 \left( 1.5147 + \frac{11.331}{K_{p.\text{пр}}} \right) + \frac{395.7}{K_{p.\text{пр}}} \\
1230.95 &= 439.26 + \frac{3286}{K_{p.\text{пр}}} + \frac{395.7}{K_{p.\text{пр}}} \\
791.69 &= \frac{3681.7}{K_{p.\text{пр}}} \implies K_{p.\text{пр}} \approx 4.65 \text{ (в разах)}
\end{align*}
В децибелах новое значение: $K_{p.\text{пр}} \text{ (дБ)} = 10 \log_{10}(4.65) \approx 6.67 \text{ дБ}$.
\textit{Вывод:} Увеличение коэффициента усиления преселектора снижает вклад шумов последующих каскадов (ПЧ и УПЧ), что существенно повышает чувствительность РПУ.

\begin{center}
    \fbox{\parbox{0.8\textwidth}{\centering \vspace{1cm} \textit{[Место для скриншота: Результат моделирования СПМ при $K_{p.\text{пр}} \approx 6.7 \text{ дБ}$]} \vspace{1cm}}}
\end{center}

\textbf{Б) Подбор коэффициента шума ПЧ ($F_{\text{пч}}$)}

Восстанавливаем $K_{p.\text{пр}} = 1.585$ (2 дБ). Найдем требуемый $F_{\text{пч}}$:
\begin{align*}
1230.95 &= 290 \left( 1.5147 + \frac{(F_{\text{пч}} - 1) \cdot 1.259}{1.585} \right) + \frac{627 \cdot 1.259}{1.585 \cdot 1.995} \\
1230.95 &= 439.26 + 230.35 (F_{\text{пч}} - 1) + 249.6 \\
542.09 &= 230.35(F_{\text{пч}} - 1) \implies F_{\text{пч}} - 1 \approx 2.353 \implies F_{\text{пч}} \approx 3.353
\end{align*}
В децибелах: $F_{\text{пч}} \text{ (дБ)} = 10 \log_{10}(3.353) \approx 5.25 \text{ дБ}$.
\textit{Вывод:} Снижение уровня собственных шумов преобразователя частоты (ПЧ) напрямую улучшает общую шумовую картину и повышает чувствительность приемника.

\begin{center}
    \fbox{\parbox{0.8\textwidth}{\centering \vspace{1cm} \textit{[Место для скриншота: Результат моделирования СПМ при $K_{\text{ш.пч}} \approx 5.25 \text{ дБ}$]} \vspace{1cm}}}
\end{center}

\textbf{В) Оценка влияния потерь в фидере}

Рассчитаем СПМ для идеального фидера ($L_{\Phi} = 0 \text{ дБ} = 1$ раз). При этом $F_{\Phi} = 1$:
\begin{align*}
F_{\text{БВЧ}}^{\text{ид}} &= 1 + (1.995 - 1)\cdot 1 + \frac{(10 - 1)\cdot 1}{1.585} = 1 + 0.995 + 5.678 \approx 7.673 \\
T_{\text{БВЧ}}^{\text{ид}} &= 290(7.673 - 1) \approx 1935.2 \text{ К} \\
K_{\text{БВЧ}}^{\text{ид}} &= 1 \cdot 1.585 \cdot 1.995 \approx 3.162 \\
T_{\text{РПУ}}^{\text{ид}} &= 1935.2 + \frac{627}{3.162} \approx 2133.5 \text{ К} \\
T_{\Sigma}^{\text{ид}} &= 300 + 2133.5 = 2433.5 \text{ К} \\
G_{ш.\Sigma}^{\text{ид}} &= 1.38 \cdot 10^{-23} \cdot 2433.5 \approx 3.358 \cdot 10^{-20} \text{ Вт/Гц}
\end{align*}
Сравнение: $\frac{3.358}{4.225} \approx 0.795$. Отсутствие потерь в фидере снижает суммарную СПМ шума на $\approx 20.5\%$.

\begin{center}
    \fbox{\parbox{0.8\textwidth}{\centering \vspace{1cm} \textit{[Место для скриншота: Результат моделирования СПМ при $L_{\Phi} = 0 \text{ дБ}$]} \vspace{1cm}}}
\end{center}

\vspace{1em}
\textbf{9. Расчет характеристик шума в различных точках приемного тракта}

\textbf{Выход преселектора:}
Коэффициент передачи от входа до выхода преселектора: $K_1 = \frac{K_{p.\text{пр}}}{L_{\Phi}} = \frac{1.585}{1.259} \approx 1.259$.
\begin{align*}
\text{Макс. СПМ: } G_{\text{вых.пр}} &= G_{ш.\Sigma} \cdot K_1 = 4.225 \cdot 10^{-20} \cdot 1.259 \approx 5.319 \cdot 10^{-20} \text{ Вт/Гц} \\
\text{Мощность шума: } P_{\text{вых.пр}} &= G_{\text{вых.пр}} \cdot \Pi_{\text{пр}} = 5.319 \cdot 10^{-20} \cdot (3 \cdot 10^6) \approx 1.596 \cdot 10^{-13} \text{ Вт}
\end{align*}

\textbf{Выход УПЧ:}
Общий коэффициент усиления тракта: $K_{\Sigma} = K_{\text{БВЧ}} \cdot K_{0.\text{упч}} = 2.512 \cdot 3000 = 7536$.
\begin{align*}
\text{Макс. СПМ дисперсии: } G_{\text{вых.упч}} &= G_{ш.\Sigma} \cdot K_{\Sigma} = 4.225 \cdot 10^{-20} \cdot 7536 \approx 3.184 \cdot 10^{-16} \text{ Вт/Гц} \\
\text{Мощность шума: } P_{\text{вых.упч}} &= G_{\text{вых.упч}} \cdot \Pi_{\text{упч}} = 3.184 \cdot 10^{-16} \cdot (1 \cdot 10^6) = 3.184 \cdot 10^{-10} \text{ Вт}
\end{align*}

Эффективное напряжение шума на выходе УПЧ (считая нагрузку $R = 50 \text{ Ом}$):
\begin{align*}
U_{\text{ш.упч}} &= \sqrt{P_{\text{вых.упч}} \cdot R} = \sqrt{3.184 \cdot 10^{-10} \cdot 50} \approx \sqrt{159.2 \cdot 10^{-10}} \approx 1.26 \cdot 10^{-4} \text{ В}
\end{align*}

\begin{center}
    \fbox{\parbox{0.8\textwidth}{\centering \vspace{1cm} \textit{[Место для скриншота: Графики спектра и эпюры шума на выходе преселектора и УПЧ]} \vspace{1cm}}}
\end{center}

\vspace{1em}
\textbf{10. Характеристики шума на выходе ФСС (Раздел 2)}

\textit{Примечание: Для моделирования ФСС (фильтра сосредоточенной селекции) подставьте значения $C$ и $L$ в соответствии с вашим вариантом из таблицы лабораторной работы. Ниже оставлено место под эпюры, получаемые после ввода расчетных номиналов.}

\begin{center}
    \fbox{\parbox{0.8\textwidth}{\centering \vspace{1cm} \textit{[Место для скриншота: АЧХ синтезированного ФСС]} \vspace{1cm}}}
\end{center}

\begin{center}
    \fbox{\parbox{0.8\textwidth}{\centering \vspace{1cm} \textit{[Место для скриншота: Гистограмма распределения вероятностей (с границами $\pm 3 U_{\text{ш}}$) и АКФ шума]} \vspace{1cm}}}
\end{center}

\section{Результаты измерений}
Данные моделирования сведены в таблицу~\ref{tab:results}.

\begin{table}[h!]
    \caption{Экспериментальные параметры схемы}
    \label{tab:results}
    \centering
    \begin{tabular}{|c|c|c|c|}
        \hline
        \textbf{№ опыта} & \textbf{Напряжение $U$, В} & \textbf{Ток $I$, мА} & \textbf{Мощность $P$, мВт} \\ \hline
        1 & 1.20 & 10.5 & 12.6 \\ \hline
        2 & 2.50 & 22.1 & 55.25 \\ \hline
        3 & 3.30 & 30.0 & 99.0 \\ \hline
    \end{tabular}
\end{table}
